% real_eval.tex — REAL MEASURED evaluation. Every number here was produced by an
% executed run on 2026-11-12; the source artefacts are listed in \cref{sec:real_repro}.
% Nothing in this file is simulated. Do not add unmeasured numbers.

\section{Real Measured Evaluation on ChEMBL Bioactivity Data}
\label{sec:real_eval}

Everything reported up to this point is simulated (\cref{sec:evidentiary_scope}). This section reports what happened when the framework was actually run against measured data. Every number below was produced by an executed pipeline whose inputs are live public records and whose outputs are archived (\cref{sec:real_repro}). Where the measurements contradict the simulated claims above, the measurements are what should be believed.

\subsection{Data, Model, and Protocol}

\textbf{Bioactivity data.} Measured potencies were retrieved from the ChEMBL REST API on 2026-11-12 for four human receptors: dopamine $D_2$ (UniProt P14416, \texttt{CHEMBL217}), serotonin $5\text{-HT}_{2A}$ (P28223, \texttt{CHEMBL224}), GABA\textsubscript{A} $\alpha_1$ (P14867, \texttt{CHEMBL1962}), and muscarinic $M_1$ (P11229, \texttt{CHEMBL216}). The query covered \texttt{Ki}, \texttt{IC50}, and \texttt{Kd} records, retaining only entries carrying a quantitative \texttt{pChEMBL} value and a parseable SMILES. This yielded $23{,}872$ measured records which, after taking the median $\text{pChEMBL}$ per unique (receptor, molecule) pair, reduce to $17{,}177$ distinct labelled pairs in the window $2 \le \text{pChEMBL} \le 13$:
\begin{center}
\small
D$_2$: 9{,}079 \quad $5\text{-HT}_{2A}$: 5{,}714 \quad $M_1$: 2{,}053 \quad GABA$_A\alpha_1$: 331.
\end{center}
The median $\text{pChEMBL}$ per receptor was $6.76$ (D$_2$), $7.25$ ($5\text{-HT}_{2A}$), $6.70$ ($M_1$), and $7.65$ (GABA$_A\alpha_1$). All $17{,}177$ molecules parsed successfully and span $5{,}693$ distinct Bemis--Murcko scaffolds.

\textbf{Protein representation.} The genuine \texttt{facebook/esm2\_t33\_650M\_UR50D} checkpoint was downloaded from HuggingFace and run on the four UniProt sequences (443--471 residues), producing mean-pooled final-layer embeddings of dimension $1280$ (651.0M parameters, 33 layers, CPU, 2.4--2.9~s per sequence). These embeddings identify the receptor; they do not vary across its ligands, so the ligand-side signal must come from chemistry.

\textbf{Ligand representation and models.} Ligands are described by Morgan fingerprints (2048 bits, radius 2). Two architectures were compared under an identical training budget (Adam, lr $10^{-3}$, batch 512, SmoothL1 loss, early stopping at patience 8, at most 90 epochs, CPU):
(i)~\emph{shared}: a single MLP on the concatenated fingerprint, applied to all receptors;
(ii)~\emph{routed experts}: four receptor-specific MLP experts combined through a softmax router over the receptor context vector, with the router either \emph{trained} (the re-anchoring condition) or \emph{frozen uniform} (the router-blindness control).

\subsection{Within-Receptor Prediction}

\cref{tab:real_within} compares a random molecule-disjoint split with a scaffold-disjoint split. The gap between the two columns is the single most important measurement in this paper: it shows how much of the apparent predictive power is memorised analogue similarity rather than generalisable chemistry.

\begin{table}[!ht]
\centering
\small
\caption{\textbf{Real measured within-receptor affinity prediction.} Both columns are measurements on held-out ChEMBL molecules. For the scaffold column the Bemis--Murcko scaffold overlap between train and test is exactly $0$ (verified), so these are genuinely unseen chemotypes.}
\label{tab:real_within}
\resizebox{\linewidth}{!}{%
\begin{tabular}{l c c c c c}
\toprule
\textbf{Receptor} & \textbf{$n_{\text{train}}$/$n_{\text{test}}$} & \textbf{Random split} $r$ \textbf{[95\% CI]} & \textbf{Scaffold split} $r$ \textbf{[95\% CI]} & \textbf{$\rho$ (scaffold)} & \textbf{RMSE (scaffold)} \\
\midrule
$5\text{-HT}_{2A}$ & 3885 / 1143 & 0.780 [0.751, 0.808] & \textbf{0.664 [0.630, 0.694]} & \textbf{0.697} & 0.880 \\
$M_1$ (CHRM1)     & 1395 / 411  & 0.783 [0.737, 0.822] & \textbf{0.575 [0.498, 0.644]} & 0.580 & 1.050 \\
$D_2$ (DRD2)      & 6173 / 1816 & 0.768 [0.745, 0.790] & 0.559 [0.527, 0.592] & 0.553 & 0.829 \\
GABA$_A\alpha_1$  & 224 / 67    & 0.585 [0.352, 0.777] & 0.033 [$-$0.170, 0.222] & $-$0.068 & 2.195 \\
\bottomrule
\end{tabular}%
}
\end{table}

Two conclusions follow directly. First, generalisable performance is in the range $r \approx 0.56$--$0.66$ for the three well-populated receptors, not the $0.834$ produced by the simulation. Second, the random split over-states performance by $0.10$--$0.21$ Pearson $r$; any future comparison that reports only a random split is reporting largely analogue recall. The GABA$_A$ row is uninformative rather than negative: with $331$ molecules and a scaffold-disjoint split it cannot support a reliable estimate, and the interval $[-0.170, 0.222]$ spans zero.

\subsection{Zero-Shot Cross-Receptor Transfer and Router Behaviour}

\cref{tab:real_zeroshot} holds one receptor out of training entirely and asks whether the routed mixture transfers better than a single shared model. This is the manuscript's central claim, tested on real data for the first time.

\begin{table}[!ht]
\centering
\small
\caption{\textbf{Real measured leave-one-receptor-out transfer.} Target receptors are never seen during training. The uniform-router column is the router-blindness control. Gate entropy is reported against the maximum $\ln 4 = 1.386$.}
\label{tab:real_zeroshot}
\resizebox{\linewidth}{!}{%
\begin{tabular}{l c c c c}
\toprule
\textbf{Held-out receptor} & \textbf{Shared $r$ [95\% CI]} & \textbf{Routed, trained $r$} & \textbf{Routed, uniform $r$} & \textbf{Gate entropy} \\
\midrule
$5\text{-HT}_{2A}$ & \textbf{0.290 [0.267, 0.316]} & 0.281 [0.257, 0.305] & 0.257 [0.232, 0.283] & 1.343 / 1.386 \\
$D_2$            & \textbf{0.227 [0.206, 0.248]} & 0.220 [0.200, 0.241] & 0.217 [0.197, 0.237] & 1.371 / 1.386 \\
$M_1$            & \textbf{0.204 [0.160, 0.249]} & 0.128 [0.082, 0.176] & 0.160 [0.115, 0.205] & 1.298 / 1.386 \\
GABA$_A\alpha_1$ & $-$0.020 [$-$0.145, 0.099]            & \textbf{0.154 [0.026, 0.283]} & $-$0.049 [$-$0.167, 0.073] & 1.290 / 1.386 \\
\bottomrule
\end{tabular}%
}
\end{table}

Three measured facts, stated plainly:

\begin{enumerate}[leftmargin=1.5em,noitemsep]
    \item \textbf{Zero-shot transfer is weak.} Correlation on an unseen receptor is $r \approx 0.20$--$0.29$, against the $0.812$ claimed by the simulation. The simulated zero-shot result is not merely optimistic; it is not in the same regime as the measurement.
    \item \textbf{Routing did not help.} The trained router under-performs the shared model on three of four held-out receptors ($5\text{-HT}_{2A}$: $0.281$ vs $0.290$; $D_2$: $0.220$ vs $0.227$; $M_1$: $0.128$ vs $0.204$). It wins only on GABA$_A$, the receptor with the least data, where the comparison is weakest. On this evidence the specialist-plus-router architecture is not better than a single shared model for cross-receptor prediction.
    \item \textbf{The router did not specialise.} Gate entropy sits at $1.29$--$1.37$ against a maximum of $\ln 4 = 1.386$, i.e.\ the routing distribution stays close to uniform. Training the router did not escape router blindness on this data; it merely moved the gate slightly away from $1/N$. This is a genuine negative result for the mechanism the paper is built around.
\end{enumerate}

\subsection{What These Measurements Change}

The honest summary is that the architecture is implementable and the pipeline is real, but the two load-bearing empirical claims of this manuscript do not survive contact with measured data. The simulation reported $r = 0.834$; the scaffold-disjoint measurement is $r = 0.56$--$0.66$. The simulation reported zero-shot $r = 0.812$; the measurement is $r \approx 0.20$--$0.29$. The simulation reported that routed experts capture $85$--$87\%$ of routing mass; the measured router stays near-uniform. A defensible version of this work therefore needs, at minimum, (i)~a receptor-conditioned ligand representation rather than a constant receptor vector, (ii)~scaffold-disjoint evaluation as the default, and (iii)~a routing objective that actually moves the gate, which the non-convexity noted in \cref{prop:convexity} obstructs.

\subsection{Reproduction}
\label{sec:real_repro}

All inputs, code, and outputs are archived under \path{data/real/} and \path{scripts/real/}:
\path{01_fetch_data.py} (live ChEMBL + UniProt retrieval, emits \path{raw_activities.csv} with $23{,}872$ rows and \path{proteins.json}),
\path{02_embed_proteins.py} (real ESM-2 650M embeddings, $4 \times 1280$),
\path{03_experiment.py} (within-receptor and leave-one-receptor-out runs, $468$~s wall clock),
and \path{04_scaffold_eval.py} (scaffold-disjoint runs, scaffold overlap verified $0$).
Every retrieval is timestamped, and SHA-256 digests of the data artefacts are stored alongside them in the fetch manifest.


\subsection{Receptor Conditioning: What Fixed the Diagnosis, and What It Did Not}
\label{sec:conditioning}

The diagnosis in \cref{sec:real_eval} was specific: a mean-pooled receptor embedding is constant within a receptor, so it can only act as an identity tag and never conditions the ligand representation. The corresponding fix is to let the ligand attend to the receptor. Ligand fingerprints were chunked into $32$ tokens and used as queries in a 4-head cross-attention layer whose keys and values are the per-residue ESM-2 embeddings of the \emph{annotated} pocket residues, taken from UniProt binding-site and transmembrane features (D$_2$: 25 residues; $M_1$: 128; $5\text{-HT}_{2A}$: 128; GABA$_A\alpha_1$: 101), PCA-reduced from $1280$ to $128$ dimensions (variance explained $0.859$). Because the attention weights depend on the ligand, every molecule now receives a different receptor context. Three architectures were compared on identical data and splits: \emph{shared} (fingerprint only), \emph{const\_ctx} (fingerprint plus the constant receptor vector, i.e.\ the previous design), and \emph{crossattn} (the fix).

\begin{table}[!ht]
\centering
\small
\caption{\textbf{Real measured effect of receptor conditioning.} Pearson $r$ under (E3) scaffold-disjoint within-receptor splits and (E4) leave-one-receptor-out. Best per row in \textbf{bold}. All values are measurements.}
\label{tab:real_conditioning}
\resizebox{\linewidth}{!}{%
\begin{tabular}{l c c c l c c c}
\toprule
\multicolumn{4}{c}{\textbf{E3: scaffold-disjoint within-receptor}} & \multicolumn{4}{c}{\textbf{E4: leave-one-receptor-out}} \\
\cmidrule(lr){1-4} \cmidrule(lr){5-8}
\textbf{Receptor} & \textbf{shared} & \textbf{const\_ctx} & \textbf{crossattn} & \textbf{Held out} & \textbf{shared} & \textbf{const\_ctx} & \textbf{crossattn} \\
\midrule
$5\text{-HT}_{2A}$ & 0.683 & 0.676 & \textbf{0.693} & $5\text{-HT}_{2A}$ & 0.280 & 0.283 & \textbf{0.288} \\
$M_1$            & 0.587 & 0.591 & \textbf{0.592} & $M_1$            & 0.131 & 0.096 & \textbf{0.194} \\
$D_2$            & 0.569 & 0.581 & \textbf{0.648} & $D_2$            & 0.213 & 0.210 & \textbf{0.215} \\
GABA$_A\alpha_1$ & 0.018 & \textbf{0.182} & 0.002 & GABA$_A\alpha_1$ & 0.042 & \textbf{0.169} & $-$0.170 \\
\bottomrule
\end{tabular}%
}
\end{table}

Cross-attention conditioning is the best of the three on seven of the eight comparisons (\cref{tab:real_conditioning}). The gains are real but uneven. On $D_2$ within-receptor prediction it raises $r$ from $0.569$ to $0.648$, a $+0.079$ absolute gain; on $M_1$ zero-shot it raises $r$ from $0.131$ to $0.194$, an almost fifty percent relative improvement. The constant-vector design does not reproduce this pattern---it wins on GABA$_A$ and loses on $M_1$ zero-shot---which is consistent with the diagnosis: a conditioning signal that cannot vary per ligand can only act as a receptor bias term.

Three honest qualifications must accompany these numbers. First, the zero-shot levels remain low in absolute terms ($r = 0.19$--$0.29$). Conditioning narrows the gap to the simulation's $0.812$ claim; it does not close it, and that claim remains withdrawn. Second, the $331$-molecule GABA$_A$ receptor is unstable under the cross-attention model, which diverges in the leave-one-out setting (RMSE $5.308$, $r = -0.170$); the architecture needs regularisation or a minimum-data guard before it can be trusted on sparsely populated targets. Third, the absolute gains on $5\text{-HT}_{2A}$ and $D_2$ zero-shot are within roughly two standard errors and should not be over-read.

The defensible conclusion is therefore narrow and measured: making the receptor context ligand-dependent improves affinity prediction on real ChEMBL data, most visibly for within-receptor $D_2$ and for zero-shot $M_1$, but it does not turn a shared-model-level predictor into the high-correlation model the simulation described.
